\documentclass[../main.tex]{subfiles}
\begin{document}
%==========================================TIÊU ĐỀ TẬP========================================
\begin{table}[h]
  \begin{adjustwidth}{-1cm}{}
    \begin{tabular}{>{\centering\arraybackslash}p{6cm}|>{\raggedright\arraybackslash}p{12.5cm}}
      \multirow{5}{6cm}
      % Cột bên trái: ÉPISODE và số tập
      {\\[-40pt]
        \flaregothic\fontsize{20pt}{20pt}\selectfont \centering
        \color{eptype}ÉPISODE\color{black}\\
        \burbank\fontsize{72pt}{72pt}\selectfont
        \color{epnum}116
        \\[18pt]
        % \flaregothic\fontsize{14pt}{14pt}\selectfont
        % \color{epattr}BIENVENUE, \uline{\color{Banpire} Mel} !
        \color{black}
      }
       &
      % Dòng thứ nhất: tên tập tiếng Nhật
      {\fotskip\fontsize{18pt}{18pt}\selectfont \ruby{変}{へん}なこと\ruby{言}{い}っちゃうんだけど！
      }                                          \\[6pt]
      % Dòng thứ hai: tên tập tiếng Anh
       & {\cafeteria\fontsize{18pt}{18pt}\selectfont Verbal Pek - The Peko Note}                                            \\[4pt]
      % Dòng thứ ba: tên tập tiếng Việt
       & {\vnmsans\fontsize{18pt}{18pt}\selectfont Cuốn sổ Peko}                                                            \\[8pt]
      % Dòng thứ tư: Nhân vật xuất hiện
       & {\sen{\fontsize{14pt}{14pt}\selectfont Track used: Delicacy Temptation (\textnormal{グルメの\ruby{誘}{ゆう}\ruby{惑}{わく}})}
         }
      \\[8pt]
      % Dòng cuối cùng: Thời gian diễn ra của tập (thường sẽ là ngày phát hành)
       & {\continuum\fontsize{16pt}{16pt}\selectfont Ngày phát hành: 01/08/2021}                                            \\[8pt]
       & (dựa theo bản dịch của \href{https://www.facebook.com/mamisubteam}{MamiSubs - Mami Fansub})                        \\[18pt]
    \end{tabular}
  \end{adjustwidth}
\end{table}
%=========================================TRUYỆN CHÍNH========================================
\color{black}

Ngày đầu tiên của tháng Tám.

Mùa hè đã bước sang tháng cuối cùng của nó, những ngày oi ả của tháng Bảy dường như vẫn còn đọng lại đâu đó trong không khí, nhưng buổi sáng nay, từng cơn gió nhẹ bắt đầu mang theo hương lá khô và mùi đất của tháng Tám. Ngoài kia, cây xanh đã lấm tấm vài chiếc lá vàng mỏng đầu tiên, như lời nhắc nhở rằng mùa thu sắp sửa chạm ngõ.

Tại văn phòng, bầu không khí vẫn còn lặng lẽ. Nhưng sự yên tĩnh ấy nhanh chóng bị phá vỡ bởi một câu chào quen thuộc của một thành viên mà ai ở công ty cũng đã nhẵn mặt: ``Kon-peko, kon-peko, kon-peko! \hll\ Thế hệ thứ ba, Usada Pekora đây, peko\~{}!!''

Cô nàng Thỏ vui tươi xông vào văn phòng, tay làm động tác đặc trưng trước ngực, mái tóc bạc buộc hai bên đung đưa theo từng bước nhún nhảy của cô. Gương mặt cô rạng rỡ như chưa từng biết đến khái niệm ``mệt mỏi'', nếu không muốn nói là tươi tỉnh hơn mọi khi.

Tuy nhiên, đi ngay sau lưng cô là một cậu con trai với gương mặt tỉnh bơ, đang nhìn quanh căn phòng và lẩm bẩm như báo cáo hiện trạng: ``Trong này trống hoắc à. Chưa có ai đến cả.''

Pekora dừng lại giữa chừng, tai thỏ vểnh lên: ``Ủa? Hôm nay không có ai đến sớm hả, peko?''

\img{9-7a}

Cô bước vòng quanh căn phòng, cúi người nhìn vào sau ghế bành, ló đầu vào phòng thu âm bên cạnh rồi bước trở lại, tay khoanh trước ngực:
``Lạ vậy\dots\ bình thường Mel hoặc Flare đến sớm lắm mà, peko\dots''

``Này, Peko-chan. Nhìn này.'' Giọng của cậu Tổng cất lên từ phía sau.

Pekora quay đầu lại, thấy cậu đang đứng gần bàn uống nước, chỉ tay xuống sàn. Cô bước lại gần, ánh mắt tò mò. Trên sàn nhà, giữa ánh nắng chiếu xiên qua cửa sổ, là một cuốn sổ tay màu xanh lam, viền bìa xanh nước biển nhạt. Cuốn sổ hướng mặt sau lên, lộ ra một chiếc mã QR.

Ở bìa trước của cuốn sổ là một cái tên được tô nổi bật với hai màu xanh lam-hồng: ``PEKO NOTE''.

\img{9-7b}

Nàng Thỏ chớp mắt vài lần, rồi nhặt cuốn sổ lên.

``Sổ tay Peko à?'' Kuon đọc lên, nheo mắt quan sát kỹ. ``Anh thấy trông quen lắm\dots\ Cái này là đồ merch của em hả, Pekora-chan?''

Cô nàng lắc đầu liên tục, tai thỏ cũng lắc theo: ``Em đâu có phát hành cái thứ kỳ cục như vậy đâu peko\~{}!''

Pekora lật cuốn sổ ra. Ở trang đầu tiên, được viết bằng chữ in đậm nét tròn là một dòng hướng dẫn: ``Điều hài hước sẽ xảy ra với bất kỳ ai có tên được ghi vào đây. Nếu xóa tên, người đó sẽ trở lại bình thường.''

Hai người cùng đọc lên, ngơ ngác nhìn nhau.

``Cái này\dots\ trông giống như một bản parody của D**th Note thì phải?''

``Em thấy nó giống Joke Note thì có, peko.''

Nhưng rồi cả hai cùng lật sang trang bên cạnh, chỉ để thấy một dòng hoàn toàn lạc quẻ được viết rất gọn gàng: ``Sau khi đã rán thịt bò với hành tây, cho thêm roux\footnote{Một loại nước xốt được làm từ bột mỳ và các loại chất béo (thường là bơ, dầu, mỡ\dots)} và đun nhỏ lửa lại.''

Cả hai cùng khựng lại vài giây, rồi cố kết nối câu này với đoạn ở trang trước, dù chắc chắn chúng chẳng hề tương quan nhau một chút nào.

``\dots Có gì đó sai sai ở đây.''

``Câu cuối này hình như là công thức làm cà ri thôi mà, peko?''

Kuon gật đầu chậm rãi, tay khoanh lại đầy nghi ngờ: ``Tự dưng có ông rảnh háng nào lại viết câu đó được nhỉ?''

Nàng ``Tội phạm chiến tranh'' nheo mắt, rồi nhìn quanh văn phòng một lần nữa, rồi cúi nhìn cuốn sổ vẫn còn thơm mùi giấy mới.

Gió ngoài trời khẽ thổi, làm rèm cửa rung lên nhè nhẹ.

Một cuốn sổ kỳ lạ. Một căn phòng vắng.

Trong lúc cả hai người vẫn đang cố tìm hiểu về nguồn gốc của cuốn sổ kỳ lạ đó, thì\dots

Tiếng cửa bật mở khe khẽ phía sau lưng họ.

Pekora vội vàng kéo tay Kuon, đẩy cả hai chui tọt vào sau tấm bảng trắng lớn dựng sát tường.

``Khoan, em làm gì vậy hả!?'' cậu phản ứng trong bối rối, suýt đụng đầu vào mép bảng.

``Suỵt.'' Nàng Thỏ đưa ngón trỏ lên môi, tai thỏ vểnh lên căng thẳng.

Bước vào phòng là một cô gái với mái tóc hồng uốn nhẹ, dáng đi đầy tự tin và gương mặt sáng rỡ như thể vừa thắng một trận đấu trò chơi điện tử kéo dài xuyên đêm: ``Nya-hello\~{}! Miko-senpai đến rồi đây\~{}!''

Không có ai trả lời. Căn phòng vẫn im lặng như tờ, như thể vẫn chưa có ai thật sự đặt chân tới đây.

\img{9-7c}

Nàng Vu nữ nghiêng đầu nhìn quanh: ``Ơ? Kỳ vậy ne? Chắc mình là người đến sớm nhất hôm nay rồi\~{}''

Cô nàng rảo bước vào trong, đi ngang qua chỗ hai người đang nấp mà không để ý gì. Kuon nghiêng đầu ra sau tấm bảng, thấy rõ nét mặt phấn khởi của cô nàng``ưu tú'' đang bắt đầu lục soát từng ngăn tủ để khám phá xem có ai cất đồ ăn vặt không.

Lúc này, Pekora cầm sẵn cuốn sổ cùng cây bút, ánh mắt lấp lánh như đứa trẻ sắp nghịch phá.``Để em viết thử tên chị ấy vào đây xem nào, peko.''

``Ờm\dots\ anh nghĩ là ta không nên đâu,'' Kuon nói nhỏ, nét mặt lộ vẻ đắn đo. ``Anh không chắc thứ này sẽ xảy ra điều gì, nhưng cũng đừng nghịch bừa-''

``Cứ viết đại đi! Có khi chuyện gì thú vị sẽ xảy ra thì sao, peko\~{}!''

Cậu nghe nàng Thỏ thuyết phục như vậy, cũng chỉ có thể thở dài: ``Anh nhớ rõ nó không phải là cuốn sổ chết người đâu ha\dots\ Thôi kệ, thử một cái chắc cũng không việc gì đâu nhỉ?''

Nhận được sự chấp thuận của cậu, Pekora hí hửng viết vào dòng đầu tiên trong trang trống: ``Sakura Miko''.

Cả hai lập tức quay đầu nhìn ra. Miko đúng lúc đó mở tủ lạnh: ``Ồ, trong này có bánh kem nè, peko\~{}!''

``\eng{\dots Cái gì vậy?}''

``Hửm? Nghe có vẻ lạ tai lắm nha\dots''

Miko cầm chiếc bánh lên, mỉm cười đầy mãn nguyện\dots\ rồi đột ngột dừng lại.

``\dots K-Khoan đã\dots'' cô lẩm bẩm. ``Mình đang nói cái peko gì thế này!?''

Cô bắt đầu hoảng loạn, ngửa đầu lên trần, mặt tái xanh như tàu lá chuối: ``Không\dots\ không phải lại nữa chứ!? Mình\dots\ mình bị sốc-peko rồi sao, peko!?''

\img{9-7d}

Pekora và Kuon nấp sau bảng nhìn nhau, cả hai gật gù như cùng hiểu ra điều gì.

``Ra là thế\dots\ Cái sổ này cũng na ná giống như cái cuốn sổ đen chết chóc kia,'' cậu Tổng ôm trán.

``Dường như bất cứ ai bị viết tên vào đây\dots\ sẽ nói chuyện như em đó, peko.''

Nàng Thỏ mím môi, mắt sáng rực như phát hiện một kho báu: ``Tức là\dots\ em có thể bắt bất cứ ai phải `peko' giống em á!? Nếu em viết ra một loạt cái tên thì sao nhỉ\dots''

``Chắc là cũng giống như Miko-chi thôi. Nhưng anh nghĩ nên dừng lại ở đây thì hơn. Chuyện này mà để A-san biết được thì\dots''

``Không nghe, không nghe! Em sẽ trở thành `peko-sama', vị thần mới của thế giới \hll, peko\~{}!''

``\dots Anh ngửi thấy có mùi điềm rồi đấy.''

Trước ánh mắt bất lực của Kuon, Pekora viết thêm vài cái tên khác. Nhưng rồi cô dừng lại một nhịp, ngoảnh sang cậu Tổng đang cố nhớ lại về cuốn sổ mà cô đang cầm trên tay.

``Để em thử viết tên anh vào đây xem nào, peko\~{}!''

Kuon chợt lùi lại nửa bước: ``Anh nghĩ không cần thiết đâu. Có thể nó không hoạt động với anh mà-''

Cô nàng Thỏ nghịch ngợm không đợi thêm. Cô viết xuống dòng tiếp theo: ``Hikawa Kuon''.

Rồi, cô nhìn cậu chằm chằm, hai tay chống hông, mặt hớn hở: ``Giờ thì nói thử một câu đi, peko!''

Cậu Tổng thở ra một hơi, rồi đáp: ``Anh đã bảo rồi, nó không có tác dụng với anh đâu. Em đừng ngớ ngẩn như vậy nữa.''

Không có một ``peko'' nào được thốt ra từ lời cậu.

``Hả!?'' Pekora ngạc nhiên cực độ. ``Sao lại không bị gì hết vậy peko!? Thử nói cho em câu khác xem!?''

Kuon chỉ tặc lưỡi rồi nói gọn: ``Đưa cho anh cuốn sổ. Có khi trong này còn có ràng buộc gì nữa.''

``K-Không thể nào, peko\dots'' nàng Thỏ ngỡ ngàng cực độ, khi vẫn không có một từ ``peko'' nào được cậu Tổng nói đến.

Không chờ Pekora trả lời, cậu lấy cuốn sổ về, lật nhanh các trang, rồi dừng lại ở bìa cuối. Một dòng chữ nhỏ được in ở mép dưới cùng: ``Không có tác dụng với những người không phải idol \hll.''

``\dots Chẳng trách,'' cậu chậm rãi nói, chỉ vào dòng chữ nhỏ đó cho đối phương xem.

Pekora cau mày: ``Ủa? Vậy là không phải ai cũng bị à\dots''

``Ừ, chắc vì anh là Tổng quản lý, là một holoStaff. Còn em thì đang cố cai trị cả dàn thần tượng luôn rồi.''

Cô nàng thỏ ngẫm một giây, rồi tuyên bố đầy quyết tâm: ``Dù sao thì, ít nhất em vẫn có thể thống trị cả đám idol, peko\~{}!''

``Anh thấy cái thế giới này sắp nổ tung đến nơi rồi\dots\ Mà, em thực sự nghiêm túc với chuyện này thật à?'' cậu nhướng mày.

``Tất nhiên rồi, peko!'' nàng Thỏ hí hửng đáp.

``\textit{Em ấy bị cuốn sổ kia điều khiển mất rồi\dots\ Chuyện này chắc chắn sẽ chẳng kết thúc êm đềm cho mà xem,}'' cậu ngước mắt lên, thở dài ngán ngẩm với mưu đồ bất chính đó.

\ct

Bên trong văn phòng giờ đã rộn ràng tiếng động, không phải vì mọi người đang làm việc, mà vì một ``hội chứng kỳ quặc'' đang lan rộng.

Không chỉ riêng Miko, mà các thành viên khác cũng đã có những triệu chứng giống như vậy: hầu như mỗi câu họ nói đều phải thêm ``peko'' vào cuối.

Rushia đập bàn bực tức, giọng cao vút càng trở nên gay gắt hơn: ``\textbf{Cái quái gì đang xảy ra thế này hảaa, peko!?}''

Noel ngồi thụp xuống ghế, hai tay ôm bụng, vẻ mặt tiu nghỉu như thể trời sập đến nơi: ``Tớ đang rất đói\dots\ peko!?''

Luna thì ôm đầu đầy khổ sở, trôngnhư sắp khóc: `Nàm sao để khắc pục điều nì đâyyy peko\dots''

Còn Miko, ngồi cạnh nàng Pháp sư, thì chỉ im lặng, không phải vì nàng ``ưu tú'' không muốn nói, mà vì bất cứ khi nào mở miệng ra, câu nói cũng bị chèn một tiếng ``peko'' vô lý vào cuối.

\begin{figure}[H]
  \centering
  \includegraphics[width=12.5cm]{../../../../Image/Book?/9-7e.png}                    % Tự động thêm "Image/" vào trước tên tệp; chỉ cần chèn tên tệp
  \caption{Rõ ràng là họ đang cùng mắc phải một thứ gọi là ``Hội chứng Peko hóa cấp tính.''\\(Dịch nghĩa dòng chữ ở hình: ``Hội những người nói peko'')}
\end{figure}

Đứng nấp sau chiếc kệ tủ phía cuối phòng, nơi đang là điểm quan sát bí mật, Pekora và Kuon đang chứng kiến cảnh tượng hỗn loạn đó, với tâm trạng hoàn toàn trái ngược.

Nàng Thỏ cầm cuốn sổ trên tay, miệng nở một nụ cười gian: ``Đế chế Peko-chan đang mở rộng phạm vi rồi, peko! Mở đầu quá thành công luôn đó, peko!''

Cậu Tổng thì khoanh tay lại, trán nhăn lại, tỏ vẻ hối hận vì sinh ra trong thời đại mà cậu gọi là ``loạn lạc'' này: ``Ta có thể dẹp cái trò chơi thần thánh này được chưa?''

Lập tức, nàng Thỏ huých vai cậu: ``Nào, đừng có làm em mất hứng! Xem kìa, họ sắp làm gì đó rồi!''

Quay trở về chiếc bàn, nàng Vu nữ hít một hơi sâu rồi vỗ tay gây chú ý tới ba người còn lại: ``Mọi người, bình-peko-tĩnh lại nào, ne!''

Rushia, Noel và Luna cùng nhìn về phía cô bằng những ánh mắt đầy tuyệt vọng, dường như rất bi quan với bản thân.

Nhưng với ánh mắt quyết liệt, cô phất tay như đang chỉ huy đội quân: ``Có khóc cũng không giải quyết được gì đâu, peko! Chúng ta phải đương đầu chấp nhận với nó và\dots\ stream ngay thôi, peko!''

Ba người kia tròn mắt nhìn cô nàng Vu nữ, bất ngờ với một giải pháp nghe\dots\ vừa vô lý vừa hợp lý một cách lạ lùng.

Miko đứng dậy, ngoảnh về phía trước, đầu khẽ ngẩng lên cao như một thủ lĩnh đích thực: ``Chị dám cá là các fan sẽ rất thích điều này, ne! Một trải nghiệm mới đầy peko! Nếu ta không thể tránh khỏi, thì hãy biến nó thành lợi thế của riêng mình, peko!''

\img{9-7f}

Noel vỗ tay hăng hái: ``Tán thành! Đây sẽ là một sự thay đổi tốt á-peko! Hơi kỳ, nhưng cũng vui!''

Rushia gật gù: ``Rushia cũng muốn thử luôn, peko! Biết đâu sẽ được donate nhiều hơn, peko-peko!''

Luna thì còn hơi chần chừ, vặn vẹo các ngón tay: ``Nhưng\dots\ đìu này lại ngược nại với chính xách `nữ zương đáng yêuuu' của em rùi đó, peko-nanora\dots''

Nàng Vu nữ nghe vậy xốc lại tinh thần: ``Không sao hết, công chúa của chị! Hôm nay, em sẽ là Nữ vương của Peko-vương quốc!''

``\dots Peko-vương quốc nghe cx hay đó, hìhì\~{}''

Và rồi, như đã tìm được tư tưởng chung, cả bốn người cùng đồng thanh: ``\textbf{Làm-peko-thôi nào!!!}''

Họ rời khỏi phòng với quyết tâm ngút trời, bước vào hành trình chinh phục khán giả bằng chính ``lời nguyền'' mới này. Dù nói năng kỳ quặc, nhưng thần thái lại vô cùng chuyên nghiệp. Cứ như thể họ đã chấp nhận sự kỳ dị ấy như một phần của bản thân, và sẵn sàng biến nó thành nội dung cho buổi phát sóng trực tiếp mới.

Đằng sau nơi ẩm nấp, hai người còn lại im lặng nhìn theo họ.

Kuon khẽ lẩm bẩm, có gì đó vừa nể phục, vừa áy náy: ``Họ thật sự\dots\ biến tai họa thành nội dung được à. Anh muốn tặng họ một tràng pháo tay thật.''

Pekora khẽ gật đầu: ``Bất ngờ thật đó, peko. Mà\dots\ em cũng thích như vậy.''

Cậu Tổng thở dài, tỏ vẻ không đồng tình với hành động điên rồ của nàng Thỏ: ``\dots Nói thật, em thấy đủ chưa? Chúng ta nên dừng trò này ở đây thì hơn.''

Pekora liếc mắt, nhếch mép cười: ``Chưa gì mà đã xoắn quầy rồi, peko? Đây mới chỉ là khởi đầu thôi\~{}''

Nói rồi, cô hí hửng lao về phía phòng thu, tay vẫn ôm khư khư cuốn sổ  như báu vật.

Kuon lững thững theo sau, chỉ biết lắc đầu, thở dài: ``Anh thấy cái này chẳng hay ho gì đâu. Mà thôi\dots\ để xem bốn người đó sẽ thích nghi thế nào\dots''

\ct

Pekora mở nắp chiếc laptop quen thuộc, gõ gõ vài phím rồi bật thẳng vào buổi stream đang diễn ra. Theo sau là Kuon, người không biết làm gì hơn ngoài việc cùng với cô nàng theo dõi buổi phát sóng của các cô gái - cũng là nạn nhân của cuốn sổ Peko kia.

``Để xem họ sẽ làm như thế nào\~{}''

``Hy vọng là mọi thứ sẽ diễn ra suôn sẻ.''

Trên màn hình là giao diện kênh của Sakura Miko, đang phát trực tiếp với tiêu đề to đùng chói lóa: ``PEKO ĐẠI CHIẾN - Livestream đặc biệt từ những nạn nhân bất đắc dĩ!''

Đi kèm là hai dòng chữ phụ đầy tính giải trí:

\begin{center}
  ``Tin sốt dẻo!''
  ``Các bạn muốn thấy Miko bị Peko-hóa làm gì?''
\end{center}

Trên màn hình, Miko xuất hiện với diện mạo quen thuộc: mái tóc hồng thướt tha, bộ trang phục vu nữ truyền thống, nhưng lần này lại được dán thêm một cặp tai thỏ rung rung trên đầu, do bộ lọc hình mà ra.

\img{9-7g1}

Cô cười toe toét trước máy quay, tay vẫy nhịp nhàng: ``Kon-peko, kon-peko, kon-peko~! Mình là \hll\ Thế hệ 0, Usada Miko-peko đây\~{}!''

Kuon ngạc nhiên khẽ bật cười: ``Anh thấy em ấy dường như\dots\ thích thú thật sự với cái trò này.''

Pekora thì toát mồ hôi hột: ``Usada' liệu có hơi quá không vậy\dots?''

Cô chưa kịp hết sốc thì nàng Vu nữ tai thỏ bắt đầu chơi mini game với khán giả, như ``Thử thách nói 10 câu mà không dính Peko!'' và ``Mỗi lần lỡ nói Peko, phải bắt chước tiếng thỏ!''.

Ngay sau đó, nàng Thỏ chuyển sang một buổi phát sóng trực tiếp chung của hai cô nàng Thế hệ thứ Ba: Noel và Rushia.

Màn hình hiện ra với nền hồng ngọt lịm, đầy hình trái tim, và hai cô gái cũng đang\dots\ cực kỳ nhập vai.

``Mình là Usada Noel-peko\~{}! Hôm nay sẽ đập giáp và đập luôn nỗi buồn, peko!''

``Còn mình là Usada Rushia-peko\~{}! Ai dám bảo chị đây không đáng yêu thì\dots\ peko-smash!!''

\img{9-7g2}

Cả hai cười đùa tung hứng qua lại, đọc bình luận fan, thi nhau nói nhanh những câu ``không peko'' rồi lập tức thất bại trong tiếng cười rộn ràng của những người xem trực tiếp qua luồng bình luận.

Pekora chớp mắt: ``Ơ\dots\ Họ chơi vui thật kìa\dots''

Kuon thì mỉm cười nhìn hai cô nàng đang vui vẻ đùa giỡn nhau đó: ``Quả thật hai đứa ấy chuyên nghiệp thật đấy.''

Không dừng lại ở đó, nàng Thỏ bật tiếp kênh của Luna, và lần này cô còn thấy một thứ gì đó động trời hơn.

Màn hình hiện lên phông nền trắng-xám u ám, với các dòng chữ đỏ viết nguệch ngoạc như ``peko\dots\ peko\dots\ không thoát được đâu\dots'' cùng những nét chữ nguệch ngoạc khác.

Luna xuất hiện với ánh sáng lờ mờ, tay cầm chiếc micro đồ chơi: ``Tớ là Usada Luna-peko\dots\ Một công chúa đã mất đi ngai zàng, chỉ còn lại tai thỏ-nanopeko\dots''

\img{9-7g3}

Kuon bật cười khẽ: ``Em ấy làm dựng cả bối cảnh kinh dị luôn kìa\dots''

Pekora trợn mắt ngồi im như pho tượng như thể không tin được vào những gì đang xảy ra: ``E-Em không biết phải cảm thấy thế nào về điều này nữa, peko\dots''

Cậu khẽ nghiêng đầu: ``Xem chừng họ không những không trốn tránh, mà còn tận hưởng nữa. Chẳng biết là bài `peko' của em còn hiệu quả không nữa.''

Pekora mở phần bình luận từ ba buổi stream để xem fan đang nói gì. Một tràng tin nhắn đổ xuống như lũ:

\begin{itemize}
  \item ``NAI-XỪ\~{} Usada Miko đáng yêu quá đi mất!!''
  \item ``Luna chơi horror vibe peko làm tớ rợn tóc gáy mà vẫn cười xỉu\~{}''
  \item ``Noel Rushia đúng kiểu chị em thỏ siêu năng động!!''
\end{itemize}
Dĩ nhiên, cũng có vài bình luận phản ứng ngược:

\begin{itemize}
  \item ``Bắt chước Pekora-sama như thế này thật không phải!''
  \item ``Trông không hợp tí nào cả\dots\ Làm ơn dừng lại đi\dots''
\end{itemize}

Nhưng số đó quá ít ỏi và nhanh chóng bị chìm nghỉm trong biển bình luận tích cực.

\img{9-7h}

Cậu Tổng nghiêng đầu, đọc những dòng bình luận cứ thế trôi dần lên trên: ``Fan phản hồi tốt ngoài mong đợi luôn.''

Pekora cắn môi, mắt vẫn không rời khỏi màn hình: ``\dots Em không biết nên tự hào hay thấy tội lỗi nữa\dots''

Cậu đặt lon nước xuống bàn, khẽ mỉm cười: ``Cái hay của livestream là vậy đó. Biến rắc rối thành trò vui, biến tai nạn thành thương hiệu. Anh thấy mấy đứa đang làm khá tốt công việc này ấy chứ.''

Nàng Thỏ thở dài một tiếng, không phải vì tiếc nuối, mà vì không biết liệu kế hoạch của cô có còn như kỳ vọng không. ``Thôi thì\dots\ cứ để xem họ đi được tới đâu, peko.''

Cho đến khi một dòng bình luận hiện ra giữa biển tin nhắn, khiến cả Pekora và Kuon phải sững người:

\begin{itemize}
  \item ``Thật lòng mà nói\dots\ mấy bản Peko này còn xịn hơn bản gốc nữa.''
\end{itemize}

Câu chữ ấy đập thẳng vào mắt Pekora như một cú tát.

``\textbf{Xịn hơn bản gốc ư!?}'' cô hét lên, mặt tái mét như vừa thấy linh hồn tổ tiên mình hiện hồn về.

Kuon nheo mắt, đọc lại dòng bình luận đó một lần nữa, rồi khoanh tay gật gù: ``Nếu nhìn khách quan thì\dots\ đúng là thế thật.''

``\textbf{Cả anh cũng cho là vậy sao!?}'' nàng Thỏ trố mắt, cảm tưởng như bị phản bội. ``Quá đáng\dots''

Cậu Tổng vẫn giữ giọng điềm tĩnh, nhưng ánh mắt cậu thì nghiêm lại rõ rệt: ``Và đó chính là điều mà anh thấy lo đấy.''

Bị động trước câu nói đầy ẩn ý của cậu, cô lắp bắp: ``Ý anh là sao chứ, peko- Khoan\dots''

Câu nói ấy như một chiếc móc câu, đưa cô nàng Thỏ đang cay cú ấy quay về thực tại với một vực sâu đáng sợ.

Mọi suy nghĩ trong đầu cô bỗng chốc rối loạn. Một viễn cảnh kinh hoàng hiện lên:

Một kênh không còn tiếng cười đặc trưng của cô do không còn sự nguyên bản.

Một cộng đồng người hâm mộ quay lưng, chạy theo những phiên bản ``Usada'' khác thay thế.

Một bảng thống kê tụt dốc không phanh, khi lượt đăng ký rơi về số 0.

Ký ức về lời cảnh báo từ cậu Tổng lạnh lùng bỗng ùa về như tiếng chuông báo tử khi lần đầu cô được nhận việc cùng với các thành viên Thế hệ thứ Ba: ``Nếu em không có một cái gì đó đặc biệt làm nên bản thân, thì coi như em chim cút rồi đấy. Ngành công nghiệp idol ở đây cạnh tranh khốc liệt lắm, chỉ có những ai khai thác được điểm mạnh của mình mới có thể trụ vững ở đây được.''

Pekora ôm đầu hoảng loạn: ``\textbf{Không được! Nếu như thế thì\dots\ thảm họa mất thôi, peko!!}''

Kuon bật cười khì: ``Ủa? Anh tưởng em muốn làm bá chủ thế giới với những người theo chủ nghĩa `peko' mà?''

``Nhưng nếu vậy thì ai còn thèm nghe em nói nữa chứ!?'' nàng Thỏ thốt lên bật lại cậu. ``Em phải hủy kế hoạch này ngay lập tức, peko!!''

Cô đẩy chiếc laptop sang một bên, lao tới bàn bên cạnh, giật lấy cục gôm, đôi mắt rực lửa quyết tâm. Tay cô nắm chặt cây tẩy như một vũ khí của chính nghĩa.

``\textbf{`Peko' là của em! Chỉ có Pekora này mới có thôi!!}''

Và rồi, cô điên cuồng tẩy xóa từng cái tên trong cuốn sổ ``ảo-ma-pekora'' đầy ám ảnh kia, miệng lẩm bẩm như niệm chú.

\img{9-7i}

Kuon khoanh tay đứng nhìn, nở nụ cười bất lực nhưng có phần nhẹ nhõm: ``Giá như em nhận ra điều đó sớm hơn\dots\ Nhưng thôi, thà muộn còn hơn không.''

Ngay khi nét chì của từng cái tên mờ dần, những hiệu ứng ``peko'' trên sóng trực tiếp cũng lần lượt biến mất. Cả bốn cô gái dừng lại, ngơ ngác trước sự thay đổi đột ngột khi cuối cùng họ cũng có thể nói được bình thường.

``Ta nên chơi trò gì tiếp theo đây ne\~{}? \dots Hể!?'' Miko sững người khi cô không còn nói được từ ``peko'' đó nữa.

Pekora vẫn chưa hề dừng tay, hét vọng về phía màn hình: ``\textbf{Trò chơi `peko' đến đây là hết rồi đó!}''

Cả Luna cũng nở một nụ cười tươi khi cuối cùng cô cũng thoát khỏi ``lời nguyền'' bí ẩn này: ``Thực tự tớ hông hỉu đượ- Hở? Tớ nói\dots\ b-thường lại được rồi-nora!?''

Rushia và Noel cũng vậy.

``Bọn mình\dots\ có thể nói chuyện bình thường lại rồi này, nanodesu!''

``Đúng thật ha! Cảm giác như vừa tỉnh khỏi giấc mơ kỳ lạ vậy đó!''

Pekora ngã vật xuống bàn, thở dốc như vừa chạy ma-ra-tông, đôi tai thỏ nhúm lại. Cậu Tổng quản lý bước tới, khoanh tay nhìn cảnh tượng trước mắt rồi nhếch môi: ``Đấy, ai bảo ham hố làm vua cơ. Rước họa vào thân thấy chưa?''

\img{9-7j}

Câu trả lời từ dưới bàn làm việc vang lên yếu ớt: ``Anh\dots\ im miệng\dots\ đi\dots''

Cô thở dốc thêm một lúc, cố gắng lấy lại nhịp hơi thở của mình, rồi thều thào: ``Giờ\dots\ vậy là\dots\ danh tính của mình\dots\ đã được bảo vệ\dots\ peko\dots''

Kuon mỉm cười: ``May cho em đấy.''

``Đừng có chọc em nữa, peko\dots''

Trong lúc cô thở phào nhẹ nhõm, cuốn sổ tai họa vô tình bị đẩy rơi khỏi mép bàn. Nó rơi xuống đất với một bụp nhẹ, rồi tự động lật đến một trang cụ thể, như có một sức mạnh vô hình điều khiển.

Trang giấy ấy hiện lên một cái tên duy nhất: ``Usada Pekora.''

Cả hai người im bặt. Không ai nói gì. Chỉ còn lại tiếng tim đập mạnh trong lồng ngực cô nàng Thỏ khi nhìn thấy tên mình trong đó.

\img{9-7k}

Cô cúi xuống, run rẩy nhặt cuốn sổ lên. Đôi mắt mở to không chớp, nhìn chằm chằm vào dòng chữ ấy.

``Không thể nào\dots\ Cái này\dots\ là sao\dots?''

Một ý nghĩ lóe lên như tia sét đánh ngang tai: ``Pekora như thế này\dots\ là do ai đó sao!?''

Đó chính là câu hỏi lớn nhất mà nàng Thỏ này đang thắc mắc.

Một cái tên được ghi sẵn trong cuốn sổ ma quái đó, lại còn được viết trên chính bìa có cách đọc giống như cô nữa.

Rốt cuộc, ai là người đã viết tên cô lên đó?
%==========================================TRUYỆN PHỤ=========================================
\omk

Kuon nhìn theo, trầm ngâm. Rồi bỗng chốc, như có một dòng ký ức tuôn trào, cậu khựng lại.

``\dots À, hình như anh nhớ ra rồi.''

``Là sao, peko?'' cô dỏng đôi tai thỏ lên.

``Anh biết tại sao lại có cuốn sổ Peko này rồi.''

``Ý anh là sao chứ?'' cô nhíu mày.

Kuon thở dài một tiếng, rồi cười trừ: ``Chuyện là thế này\dots''

\begin{center}
  \uline{Hồi tưởng\\Đầu tháng 12 năm 2019.}
\end{center}

Một chiều muộn trong văn phòng \hll. Đèn sáng trắng hắt xuống bàn họp nơi Kuon đang ngồi cùng Miharu, Tokue - hai người bạn học cùng đại học - và A-chan. Trên bàn là một tập hồ sơ dày với các bản lý lịch mới toanh của năm ứng viên từ Thế hệ thứ Ba.

``Vậy\dots\ cơ bản là Gen 3 này chúng ta sẽ có một Nữ Hiệp sĩ, một Cô đồng, một Yêu tinh lai người, một cosplayer Thuyền trưởng, và một\dots\ con Thỏ à?'' cậu Tổng liếc về năm tờ sơ yếu, nhíu mày lại như đang soi xét rất kỹ lưỡng.

Nàng Đeo kính nhẹ nhàng gật đầu: ``Đúng rồi đó. YAGOO-sama đã chọn xong ứng viên rồi. Nhưng giờ chúng ta cần dựng tính cách cho từng người. Hikawa-san, cậu tìm hiểu đến đâu rồi?''

Cậu giở hồ hơ của từng người ra, rồi chỉ vào một trong các ứng viên được tuyển: ``Trong bản CV của cả năm người thì có lẽ Pekora này là dị hợm nhất.'' Cậu giơ một tờ lý lịch có hình minh họa một cô thỏ tóc xanh lên cho A-chan xem.

``Tôi cũng thấy thế,'' Miharu gật gù, ``Một con thỏ mà đi làm\dots\ tội phạm chiến tranh á? Mới nghe lần đầu đấy.''

``Đúng kiểu `điên nhẹ' mà ông thích ấy, Kuon.'' Tokue bông đùa.

Kuon thở dài, phì cười: ``Vấn đề của em ấy không nằm ở cốt truyện hay là ngoại hình, vì tôi không có ý kiến gì rồi. Nó lại nằm ở việc em ấy\dots\ nói chuyện nhạt quá.''

A-chan ngơ mặt ra, hỏi lại để xác nhận: ``Ý cậu là sao?''

``Tức là thế này,'' cậu chậm rãi giải thích, ``cá nhân em ấy nhìn tổng thể thì ổn, nhưng chưa có thứ gì đặc trưng khiến khán giả nhớ lâu, mà cách dễ nhận biết nhất nằm ở việc em ấy giao tiếp với khán giả như thế nào, hay dễ hiểu hơn là có cách nói như thế nào để làm nên\dots\ kiểu như, `thương hiệu' ấy.''

``Ra thế\dots'' cả ba người trong phòng họp đều hiểu rõ vấn đề. ``Quả thật như cậu nói, Usada-san thật sự chưa có một cách nào để làm nên con người cô ấy cả,'' nàng Đeo kính lập luận.

``Và vì em thấy Pekora-chan là nhân tố có tiềm năng nổi bật nhất, nên em muốn tạo cho cô ấy một cách nói chuyện thật đặc biệt.''

Cậu bạn Hiến máu rướn mày: ``Ý ông là giống kiểu Haato-chan với ``Haachama-chama\~{}'' ấy hả?''

``Chính xác luôn!'' Kuon búng tay. ``Cái đó gọi là `verbal tic'. Một tật ngôn ngữ lặp lại sau mỗi câu, khiến người ta dễ ghi nhớ và dễ nhận diện. Ví dụ như ``zura\~{}'' của Hoa Tròn trong `L*ve L*ve!' ấy.''

A-chan chỉnh lại kính, khẽ gật đầu như đã hiểu: ``À, tôi hiểu rồi. Vậy cậu định dùng từ gì?''

Cậu Tổng ngẫm một lúc, rồi không tiếc lời mà dứt khoát đáp lại ngắn gọn: ``Em dùng `peko' ạ.''

Miharu rướn mày: ``Sao Nnghe như tiếng đánh rắm vậy?''

Tokue thì bật cười nhẹ, nhưng cũng thắc mắc với lựa chọn của cậu: ``Tôi tưởng ông định nói `kora' hay `nyan' gì cơ.''

Kuon lắc đầu nhẹ: ``Nếu mà dùng `kora' thì có khác nào mắng người ta, kiểu `\textbf{Mày lại đây cho tao!}' đâu. Còn `pera' thì nghe rất buồn cười mà cũng chẳng xuôi tai tẹo nào.''

``Với cả `nyan' là dùng cho mèo mà. Cô ấy là thỏ chứ có phải là mèo đâu,'' cậu bạn Lập trình bắt bẻ.

``À, ừ, xin lỗi.''

Cậu Tổng tiếp tục giải thích: ``Vậy nên, dùng `peko' thì vừa dễ nói, vừa có chút\dots\ ngố tàu nhưng cũng có phần dễ thương.''

``Hmmm\dots'' nàng Đeo kính khoanh tay trầm ngâm với ý tưởng của cậu, rồi ra dấu đồng tình. ``Đúng là có nét riêng thật. Vậy cậu định áp dụng như thế nào? Cậu biết là ta không thể can thiệp được tâm trí vào cô ấy, đúng chứ?''

``Không phải là không thể,'' cậu nói. ``Mà là có một cách gián tiếp. Em định sẽ tạo ra một cuốn sổ đặc biệt, tựa tựa giống `D**th Note' đấy ạ. Nhưng thay vì mang tính chết chóc, nó lại giống một quyển `verbal tic note' hơn ạ. Mỗi người bị ghi tên vào sẽ bị ảnh hưởng bởi một kiểu nói chuyện đặc trưng nào đó.''

``\dots Cậu nói thế làm tôi nổi da gà đấy. Dùng công nghệ kiểu gì mà làm được?'' A-chan toát mồ hôi hột.

Không hề ngần ngại gì, Kuon chỉ về hai cậu bạn đang ngồi kế bên: ``Tất nhiên là em sẽ không làm một mình. Em còn có hai ông anh tiến sĩ máy tính tương lại đây. Một người chuyên về giao diện ngôn ngữ học AI, một người chuyên thiết kế mạch. Em còn tính gọi thêm hai người nữa cơ.''

Cậu bạn Hiến máu ra vẻ tự hào: ``Tôi sẽ lo phần mạch tích hợp phản ứng từ ngữ.''

Cậu bạn Lập trình cũng tự tin không kém: ``Còn tôi thiết kế trình kích hoạt dựa trên ngữ cảnh hội thoại.''

``Thế là\dots\ ba đứa định tạo ra một cuốn `Peko-note' thật à?''

Kuon nghiêm túc đáp lại: ``Dạ. Và nó sẽ chỉ ảnh hưởng tới cách nói, không làm thay đổi nội dung ý chí, nên không vi phạm nhân quyền.''

Sau một hồi tiếp nhận thông tin và ý tưởng từ ba kĩ sư máy tính, A-chan thở dài, rồi cười nhẹ chấp thuận: ``Thôi được. Tôi sẽ để ba đứa thử xem sao. Nhưng nhớ kỹ: đừng để lộ ra ngoài. Nếu thất bại thì coi như chưa từng có chuyện này.''

``\eng{\textbf{Rõ!!!}}'' cả ba cậu con trai giơ tay chào kiểu nhà binh.

Vài tuần sau cuộc họp với A-chan, một căn phòng nhỏ trong khu văn phòng phụ được trưng dụng làm nơi nghiên cứu bí mật. Những hộp mì ăn liền, những hộp cơm rỗng vứt đầy trong góc. Trên bàn là hàng đống linh kiện điện tử, dây nối, bảng mạch, cùng một cuốn sổ bìa trơn đang nằm giữa đống bút dạ và thước vẽ.

Ba người con trai của lớp Khoa họ Máy tính - Kuon, Miharu và Tokue - vẫn đang lúi húi dưới ánh đèn vàng nhạt. Không khí căng như dây đàn, nhưng cũng mang chút háo hức kỳ lạ.

``\vie{Chỗ mạch logic này\dots\ nếu ghi đè nhiều hơn ba tên, thì nó sẽ bắt đầu bị loạn,}'' Tokue càu nhàu, mắt vẫn dán vào mạch vi xử lý mà anh vừa hàn xong.

``\vie{Vậy thì ta giới hạn chúng lại là xong,}'' Miharu nói, vừa mở laptop kiểm tra câu lệnh vừa gõ tiếp dòng lệnh. ``\vie{Tôi sẽ viết trình xử lý dữ liệu chỉ nhận diện nếu cái tên khớp với ID trong hệ thống talent \hll. Bên ngoài thì không có hiệu lực.}''

``\vie{Ừ, có lý,}'' Kuon gật đầu. ``\vie{Ta sẽ không muốn cái sổ này ảnh hưởng lung tung tới người thường. Để phòng bất trắc, nó chỉ kích hoạt với những người có hồ sơ nội bộ lấy từ A-san thôi.}''

Một lát sau, khi phần logic đã gần như hoàn tất, Kuon đứng dậy, duỗi lưng rồi cầm điện thoại. ``\vie{Để tôi gọi Yamazu lên. Nó làm bên thiết kế cho câu lạc bộ truyền thông nên nó vẽ đẹp đấy.}''

Tokue ngẩng lên: ``\vie{Cái ông học bá cùng nhóm với ông đấy á?}''

``\vie{Chính nó. Cần một cái bìa khiến người ta cảm thấy `có gì đó kỳ lạ' từ lần đầu nhìn vào. Tôi thì có ý tưởng ban đầu rồi, để tí tôi giải thích sau.}''

Vài giờ sau, Yamazu xuất hiện với một chiếc máy tính xách tay hầm hố dùng cho việc chơi game, cùng với một vài tờ nháp giấy A4 vẽ mĩ thuật. ``\vie{Ông Kuon gọi tôi lên giúp đây này.}''

Cậu mở máy tính ra, vào phần mềm dựng ảnh kỹ thuật rồi đẩy cặp kính cận lên hỏi mọi người: ``\vie{Kuon, cho tôi ý tưởng.}''

Cậu Tổng nhíu mày, vừa giơ từng ngón tay vừa liệt kê: ``\vie{Pha chút dễ thương, nhưng phải có cảm giác như\dots\ một thứ gì đó nằm bên dưới bề mặt. Như thể nhìn vào là thấy quen, mà cũng thấy sai sai. Tôi nghĩ là cần có hai màu xanh lam và hồng, vì chúng là màu sắc chủ đạo của Pekora-chan.}''

Cậu bạn Đạo diễn không trả lời. Cậu chỉ gật đầu, và bắt đầu tô những đường màu đầu tiên, với màu xanh biển với lớp viền hồng pastel nhạt, đan xen bởi vài họa tiết xoắn ẩn chìm. Sau vài lần chỉnh sửa, bìa sổ đã có hình dáng như bây giờ: vừa đơn giản, vừa khiến người ta không thể rời mắt.

Cùng lúc đó Tamachi, một trong hai thành viên cùng nhóm ``Đồ án'' của Kuon cũng tới. Cậu được Kuon đến để trao cơ hội làm dự án này, khi cậu có một chút kiến thức về server, và giờ đang lập một danh sách mã hóa toàn bộ tên gọi nội bộ của tài năng \hll.

``\vie{Danh sách này sẽ là khoá kích hoạt,}'' Tamachi nói khi kiểm tra từng ID. ``\vie{Nếu tên viết trong sổ trùng 100\% với tên định danh trong server, sổ sẽ bắt đầu ghi lại ngữ điệu. Ngược lại thì không có tác dụng.}

``\vie{Chính xác cái tôi cần,}'' Kuon hài lòng. ``\vie{Không bõ công làm bài đồ án đa nền tảng chứ nhỉ?}''

Sau nhiều đêm không ngủ, nhóm nhỏ của họ bắt đầu chạy thử nghiệm. Lần lượt từng người trong nhóm viết tên của mình vào cuốn sổ, để kiểm tra phản ứng.`

``\vie{Tôi viết tên của ông vào rồi đấy. Ông thử nói xem có gì khác biệt không?}'' Miharu hỏi.

``\vie{Chưa có gì,}'' Tokue đáp, ``\vie{vẫn thấy bình thường, peko.}''

Nhưng rồi chính cậu bạn giật mình khi đã vô thức nói thêm một từ nữa ở cuối câu. ``\vie{Ê, ê, ê, này! Tôi nói câu này có `peko' đâu mà lại bật ra thế, peko!? Ặc, lại nữa rồi!}''

``Tốt. Vậy là máy kích hoạt đúng cách,'' cậu bạn Lập trình cười, nhập thêm vài dòng lệnh để chỉnh tốc độ phản hồi.

``\vie{Tôi chẳng biết nên vui hay lo lắng nữa\dots}'' cậu Hiến máu thở dài. ``\vie{Dọa tôi sợ chết khiếp rồi, peko- Mà này! Xóa tên tôi cái database đi cái!}''

Tuy vậy, một vấn đề lập tức được phát hiện: hệ thống phản hồi ngôn ngữ vẫn còn quá rộng, đôi khi áp dụng cả những biến thể ngẫu nhiên nghe không tự nhiên (ví dụ như `pera' hay `pero'). Cả nhóm mất thêm hai ngày nữa để tinh chỉnh bộ lọc từ khóa, tạo ra giới hạn tần suất xuất hiện, và bổ sung tính năng ``nhận diện nhân cách'',  tức là sổ sẽ chỉ áp hiệu ứng lên người có mức độ trùng khớp hành vi cao với mô tả tính cách định sẵn.

Sau tất cả, cuối cùng họ có được một cuốn sổ hoạt động như dự định: khi viết tên một thành viên thuộc tài năng \hll\ vào, nếu người đó đủ điều kiện, một verbal tic (thói quen nói) sẽ hình thành và bám theo họ như một thói quen tự nhiên.

Cái tên đầu tiên được viết vào, như một lời hứa với tương lai, không gì khác là ``Usada Pekora''.

``\vie{Bọn mình viết tên cô ấy vào nhé?}'' Kuon nói, tay cầm bút chần chừ.

``\vie{Hẳn em ấy sẽ thấy ngạc nhiên đấy,}'' Miharu nhắc, ``\vie{biết đâu nó giúp xây dựng cá tính thì cứ thử xem.}''

Một nét bút mảnh, mềm, màu xanh lam. Cái tên ``Usada Pekora'' xuất hiện, trôi nổi trên trang giấy như thể phát sáng một khắc rồi chìm vào bề mặt. Không ai biết lúc đó điều gì thực sự đã xảy ra. Nhưng ai cũng có cảm giác, một thứ gì đó đã bắt đầu.

Sau cùng, cậu Tổng cẩn thận bọc sổ lại trong vải lụa, đặt vào một ngăn kéo bí mật phía sau kệ tài liệu trong văn phòng phụ, nơi không ai, trừ họ, biết đến.

\begin{center}
  \uline{Kết thúc hồi tưởng.}
\end{center}

``Chúng ta để nó ở đây. Và quên đi,'' cậu nói.

Và họ đã thật sự quên.

Cho đến ngày hôm nay.

``Bây giờ anh mới sực nhớ ra đấy,'' cậu thở dài, tay vẫn cầm cuốn sổ đã ngả màu theo thời gian. ``Cho anh xin lỗi nha, Pekora-chan. Có lẽ là em ghét anh kể từ khi biết được thông tin chấn động này.''

Nhưng trái ngược với suy nghĩ có phần bi quan từ cậu, nàng Thỏ chỉ nhún vai, rồi bật cười khúc khích như thể chẳng có gì đáng trách. ``Thôi không sao đâu, peko! Em nghĩ nó đã giúp em trở thành một con người độc nhất vô nhị trong \hll\ đó, Kuon-senpai à.''

``Thế á?'' cậu hơi nhướng mày, có phần ngạc nhiên trước sự điềm tĩnh của cô nàng.

``Đúng á!'' cô gật đầu chắc nịch, hai tay ôm chặt lấy cuốn sổ như bảo vật. ``Nhờ vậy mà em có được cách nói chuyện không ai bắt chước được. Có riêng cho mình một con người riêng biệt hẳn, peko! Nên là cho Pekora-chan cảm ơn anh nha!'' cô hí hửng quay tròn một vòng ngay giữa phòng.

Tuy nhiên, niềm vui ấy chỉ kéo dài trong chốc lát. Pekora chợt dừng lại. Mắt cô khẽ nheo lại, như vừa nhớ ra điều gì đó. Cô lật giở vài trang đầu tiên của cuốn sổ, rồi lên tiếng, giọng nhỏ hẳn đi: ``Nhưng\dots\ có một chuyện lạ, peko.''

Cậu Tổng ngẩng lên nhìn cô, thấy rõ vẻ trầm ngâm chưa từng có trên gương mặt thường ngày luôn tràn đầy năng lượng của cô nàng. Một khuôn mặt thường thấy khi cô nàng kết thúc buổi phát sóng của mình: một cô gái nhút nhát, hướng nội và khá khép kín.

Cậu không nói một lời gì, để cho Pekora tiếp tục: ``Em là công chúa của vương quốc Pekoland, đúng không? Ở đó, `peko' là từ duy nhất trong ngôn ngữ quốc gia em. Vậy thì\dots''

Cô ngập ngừng, rồi tiếp: ``Tại sao ngay từ khi em đến Nhật, em lại không thể nói được từ đó trong tiếng người? Em đã thử đấy, peko. Nhưng mỗi lần em mở miệng, nó cứ không ra được. Dù có cố gắng thế nào, từ mẹ đẻ đó, từ độc nhất vô nhị làm nên con người em đó\dots\ lại không thể thốt ra được.''

Cả căn phòng bỗng chốc trùng xuống.

Nhưng rồi, Pekora vẫn tiếp tục bộc bạch: ``Mãi cho đến khi\dots\ vào một ngày nọ, bỗng nhiên em có thể nói được từ đó ra, peko.''

Cô nàng quay về phía cậu, khuôn mặt tỏ rõ sự cảm kích: ``Anh có biết ngay khi em thốt ra được từ đó, em đã cảm thấy thế nào không?''

Kuon im lặng, chỉ ``hừm'' một tiếng, và lắc đầu nhẹ.

``Em đã rất vui sướng đấy. Từ cúng cơm, từ mẹ đẻ mà em đã muốn thốt ra từ lâu\dots\ cuoosoi cùng cũng được bật ra. Ngày hôm đó, em mừng tới mức phát khóc luôn á, peko.''

Cô cúi đầu nhìn cuốn sổ trên tay. ``Thật kỳ lạ, peko. Có khi nào\dots\ em không thực sự đến từ Pekoland như em nghĩ? Hay là\dots\ có điều gì đó em chưa hiểu hết về chính mình?'' cô ngẩng lên trần, tỏ vẻ nghi ngờ về nhân sinh quan của bản thân, rồi hướng về phía cậu, chờ đợi một câu trả lời thỏa đáng.

Kuon im lặng. Cậu không biết nên trả lời thế nào. Đối diện với ánh mắt ấy, cái ánh mắt không còn chỉ là một tài năng vui tính nghịch ngợm và ranh mãnh, mà là một cô gái đang đứng giữa ranh giới của bản sắc và niềm tin, cậu chỉ có thể thở dài, rồi nhẹ giọng: ``Anh không biết, Peko-chan. Nhưng anh nghĩ, điều quan trọng là em cảm thấy em là ai. Nếu em tin rằng mình đến từ Pekoland, nếu `peko' là một phần khiến em hạnh phúc, một phần khiến em cảm thấy tự hào, thì\dots\ có lẽ nó đã là sự thật rồi.''

Pekora nhìn cậu, gật đầu chấp nhận với câu trả lời có phần mông lung nhưng thẳng thắn này. Cô khẽ mỉm cười, ôm chặt cuốn sổ hơn: ``Miễn là các fan yêu quý em là em vui rồi, peko!''

Như chẳng còn bận tâm tới nỗi băn khoăn vừa thoáng qua, cô nàng nhảy chân sáo ra khỏi phòng với tiếng ``peko\~{}'' vang vọng, khiến không khí cũng trở nên nhẹ nhõm hơn.

Kuon chỉ biết nhìn theo, khẽ lẩm bẩm: ``\vie{Mình thật sự\dots\ giúp em ấy sao?}''

Nhưng rồi cậu chỉ nhún vai kiểu ``biết chết liền''  rồi quay về chỗ ngồi. Cuốn sổ Peko,  giờ đây không chỉ là một công cụ thử nghiệm, mà đã trở thành một phần trong hành trình tự khám phá của một cô gái đặc biệt, được Pekora mang về nhà và cất giữ như báu vật.

Riêng về phía bốn người đã tạo ra cuốn sổ năm xưa, Kuon phải mất kha khá thời gian để trấn an họ. Khi được hỏi có nên hủy nó không, cậu đành nói dối nửa thật nửa đùa: ``Gieo nhân nào thì gặt quả ấy rồi. Pekora-chan sẽ mãi mãi nói `peko' như vậy thôi.'' (Dù thật ra, chính cô ấy đã chọn tiếp tục như thế.)

Cũng may, câu nói đó đủ để khiến Yamazu, Tamachi, Miharu và Tokue không còn lăn tăn. Không có vụ ẩu đả nào nổ ra, và ``tội phạm chiến tranh'' vẫn tiếp tục kế hoạch bá chủ thế giới- à không, bỏ ngỏ nó  hoàn toàn luôn, như thể chưa từng nhắc lại chuyện này bao giờ.

Mọi thứ rồi cũng trở lại bình thường, theo cách không ai có thể định nghĩa được thế nào là ``bình thường'' trong thế giới của \hll\ này.
\end{document}